\documentclass[12pt,a4paper]{article}
\usepackage[utf8]{vietnam}
\usepackage{amsmath,amssymb}
\usepackage{geometry}
\geometry{
  top=12mm,
  bottom=12mm,
  left=14mm,
  right=14mm,
  footskip=6mm,
  headheight=14pt
}

\usepackage{tikz}
\usetikzlibrary{arrows.meta,calc,shapes.geometric}
\usepackage[most]{tcolorbox}
\usepackage{fontawesome}
\usepackage{fancyhdr}
\usepackage{enumitem}

\pagestyle{fancy}
\fancyhf{}
\fancyhead[L]{\footnotesize\color{teal!80!black}\faBook\ \textbf{SỔ TAY KIẾN THỨC TOÁN 3}}
\fancyhead[R]{\footnotesize\color{teal!80!black}\textbf{Lớp Toán Thầy Lĩnh -- loptoanthaylinh.vn}}
\fancyfoot[C]{\footnotesize\color{gray}-- Sổ tay kiến thức: Bài 22. Luyện tập chung hình học --}
\renewcommand{\headrulewidth}{0.4pt}
\renewcommand{\footrulewidth}{0pt}

\setlength{\parindent}{0pt}
\setlength{\parskip}{0pt}

\begin{document}

% --- BANNER BÀI HỌC ---
\begin{tcolorbox}[
  enhanced,
  colback=teal!8!white,
  colframe=teal!70!black,
  arc=3mm, boxrule=1.2pt,
  left=3mm, right=3mm, top=2mm, bottom=2mm,
  halign=center
]
  {\LARGE\bfseries\color{teal!90!black}\faBookmark\ BÀI 22: LUYỆN TẬP CHUNG HÌNH HỌC}\\[1mm]
  {\itshape\footnotesize\color{black!75} Đại tổng kết hình phẳng và hình khối: Điểm, đoạn thẳng, hình tròn, góc, khối hộp trên 1 trang A4}
\end{tcolorbox}

\vspace{-2mm}

\noindent
\begin{minipage}[t]{0.49\textwidth}\vspace{0pt}
\begin{tcolorbox}[
  enhanced,
  colback=blue!4!white,
  colframe=blue!75!black,
  arc=2mm, boxrule=0.8pt,
  left=2mm, right=2mm, top=1.8mm, bottom=1.8mm,
  title={\footnotesize\bfseries\color{white}\faCheckSquare\ 1. Mạch Hình học phẳng}
]
{\footnotesize
\begin{itemize}\setlength{\itemsep}{1pt}\setlength{\leftskip}{-2.5mm}
  \item \textbf{Đoạn thẳng \& Trung điểm:} $M$ là trung điểm của $AB \implies MA = MB = AB : 2$.
  \item \textbf{Hình tròn $(O; R)$:} Đường kính $d = 2 \times R$.
  \item \textbf{Góc:} Dùng ê-ke để kiểm tra góc vuông.
  \item \textbf{Hình chữ nhật \& vuông:}
  \begin{itemize}\setlength{\itemsep}{0.5pt}\setlength{\leftskip}{-2.5mm}
    \item Hình chữ nhật: $4$ góc vuông, $2$ dài bằng nhau, $2$ rộng bằng nhau.
    \item Hình vuông: $4$ góc vuông, $4$ cạnh bằng nhau.
  \end{itemize}
\end{itemize}
}
\end{tcolorbox}
\end{minipage}
\hfill
\begin{minipage}[t]{0.49\textwidth}\vspace{0pt}
\begin{tcolorbox}[
  enhanced,
  colback=teal!4!white,
  colframe=teal!75!black,
  arc=2mm, boxrule=0.8pt,
  left=2mm, right=2mm, top=1.8mm, bottom=1.8mm,
  title={\footnotesize\bfseries\color{white}\faCube\ 2. Mạch Hình khối không gian}
]
{\footnotesize
\begin{itemize}\setlength{\itemsep}{1pt}\setlength{\leftskip}{-2.5mm}
  \item Bộ thông số bất biến: $\mathbf{6\text{ mặt}} - \mathbf{8\text{ đỉnh}} - \mathbf{12\text{ cạnh}}$.
  \item \textbf{So sánh hai khối:}
  \begin{itemize}\setlength{\itemsep}{0.5pt}\setlength{\leftskip}{-2.5mm}
    \item Khối hộp chữ nhật: $6$ mặt là hình chữ nhật.
    \item Khối lập phương: $6$ mặt là hình vuông bằng nhau, $12$ cạnh dài bằng nhau.
  \end{itemize}
  \item \textbf{Kĩ năng quan sát:} Đếm số khối nhỏ bị che khuất trong các mô hình xếp tầng.
\end{itemize}
}
\end{tcolorbox}
\end{minipage}

\vspace{-2.5mm}
% --- KHỐI 3: SƠ ĐỒ TƯ DUY TIKZ ---
\noindent
\begin{tcolorbox}[
  enhanced,
  colback=orange!4!white,
  colframe=orange!80!black,
  arc=2mm, boxrule=0.8pt,
  left=2mm, right=2mm, top=1.5mm, bottom=1.5mm,
  title={\footnotesize\bfseries\color{white}\faSitemap\ SƠ ĐỒ TƯ DUY: CÂY ĐẠI THỤ KIẾN THỨC HÌNH HỌC LỚP 3}
]
\centering
\begin{tikzpicture}[
  box/.style={draw, rounded corners=2pt, align=center, font=\scriptsize\bfseries, inner sep=2.5pt, minimum height=0.75cm},
  arr/.style={-{Stealth[length=2mm]}, thick}
]
  \node[box, fill=purple!15, draw=purple!80!black, text width=2.6cm] (h1) at (0, 0) {TRUNG ĐIỂM\\Ở giữa\\$AM = MB = \frac{AB}{2}$};
  \node[box, fill=blue!12, draw=blue!80!black, text width=2.6cm] (h2) at (3.5, 0) {HÌNH TRÒN\\Tâm $O$, Bán kính $R$\\Đường kính $d = 2R$};
  \node[box, fill=teal!12, draw=teal!80!black, text width=2.6cm] (h3) at (7.0, 0) {GÓC VUÔNG\\Kiểm tra bằng ê-ke\\Hình vuông, chữ nhật};
  \node[box, fill=orange!15, draw=orange!80!black, text width=2.6cm] (h4) at (10.5, 0) {HÌNH KHỐI\\$8$ đỉnh, $12$ cạnh\\$6$ mặt đối xứng};

  \draw[arr, purple!80!black] (h1) -- (h2);
  \draw[arr, blue!80!black] (h2) -- (h3);
  \draw[arr, teal!80!black] (h3) -- (h4);
\end{tikzpicture}
\end{tcolorbox}

\vspace{-2.5mm}
% --- KHỐI 4: VÍ DỤ & MẸO TRÁNH BẪY ---
\noindent
\begin{minipage}[t]{0.49\textwidth}\vspace{0pt}
\begin{tcolorbox}[
  enhanced,
  colback=green!4!white,
  colframe=green!60!black,
  arc=2mm, boxrule=0.8pt,
  left=2mm, right=2mm, top=1.8mm, bottom=1.8mm,
  title={\footnotesize\bfseries\color{white}\faCheckCircle\ Ví dụ thực hành chuẩn mực}
]
{\footnotesize
\textbf{1. Tính độ dài:}\\
Biết $O$ là trung điểm đoạn thẳng $AB$, $AB = 16\text{ cm}$. Bán kính của hình tròn tâm $O$ đi qua $A$ là: $R = 16 : 2 = \mathbf{8\text{ cm}}$.\\
\textbf{2. Đếm số góc vuông:}\\
Hình bên gồm $1$ hình chữ nhật ghép với $1$ hình tam giác vuông. Tổng số góc vuông là: $4 + 1 = \mathbf{5}$ góc vuông.
}
\end{tcolorbox}
\end{minipage}
\hfill
\begin{minipage}[t]{0.49\textwidth}\vspace{0pt}
\begin{tcolorbox}[
  enhanced,
  colback=red!4!white,
  colframe=red!70!black,
  arc=2mm, boxrule=0.8pt,
  left=2mm, right=2mm, top=1.8mm, bottom=1.8mm,
  title={\footnotesize\bfseries\color{white}\faExclamationTriangle\ Mẹo tránh bẫy trắc nghiệm}
]
{\footnotesize
\begin{itemize}\setlength{\itemsep}{1.5pt}\setlength{\leftskip}{-2.5mm}
  \item \textbf{Bẫy đếm thiếu khối bị che khuất:}\\
  Khi quan sát hình khối xếp tầng 3D, hãy đếm theo từng tầng từ dưới lên trên để không bỏ sót các khối lập phương ở tầng dưới đang đỡ các khối tầng trên!
  \item \textbf{Bẫy nhầm bán kính và đường kính:}\\
  Luôn đọc kĩ đề hỏi "bán kính" hay "đường kính" trước khi đặt bút tính.
\end{itemize}
}
\end{tcolorbox}
\end{minipage}


\end{document}
